% theme: strata_and_sedimentary_rocks
\MajorQuestion{地層と堆積岩}
\SetRowSpaceQanda

\Instruction{岩石がくずれ、土砂が積もるまでのはたらきに関する次の問いに答えなさい。}
\QQ{岩石が温度変化や水などの影響を受け、表面からしだいにくずれていくことを何というか。}{}
\QQ{流水が岩石や土地をけずるはたらきを何というか。}{}
\QQ{流水がけずり取った土砂を運ぶはたらきを何というか。}{}
\QQ{運ばれた土砂が水底などに積もることを\NamedBlank{deposition}という。}{}
\QQ{川によって運ばれた土砂が河口付近に\RefBlank{deposition}してできる、低く平らな地形を何というか。}{}

\Instruction{重なった層の観察と、力による変化に関する次の問いに答えなさい。}
\QQ{重なった層が地表に現れているところを何というか。}{}
\QQ{ある地点の層の上下関係や厚さ、特徴を柱のように表した図を何というか。}{}
\QQ{離れた場所の層を比べるときに目印となる、広い範囲に同じ時期に積もった特徴的な層を何というか。}{}
\QQ{横から強い力を受けて、波を打つように曲がった層のつくりを何というか。}{}
\QQ{大地にはたらく強い力によって、層が切れてずれたものを何というか。}{}

\Instruction{粒の大きさや、もとになった物質による岩石の分類に関する次の問いに答えなさい。}
\QQ{水底などに積もった物質が、上からの重みで押し固められてできた岩石を\NamedBlank{sedimentary_rock}という。}{}
\QQ{\RefBlank{sedimentary_rock}のうち、直径$2$\,mm以上の粒が押し固められてできたものを何というか。}{}
\QQ{\RefBlank{sedimentary_rock}のうち、直径$\dfrac{1}{16}$\,mm以上$2$\,mm未満の粒が押し固められてできたものを何というか。}{}
\QQ{\RefBlank{sedimentary_rock}のうち、おもに直径$\dfrac{1}{16}$\,mm未満の粒が押し固められてできたものを何というか。}{}
\QQ{\RefBlank{sedimentary_rock}のうち、生物の死がいなどが積もってでき、うすい塩酸をかけると二酸化炭素が発生するものを何というか。}{}

\Instruction{生物や火山のはたらきに由来する岩石と、過去の生物に関する次の問いに答えなさい。}
\QQ{\RefBlank{sedimentary_rock}のうち、生物の死がいなどが積もってでき、うすい塩酸をかけても気体が発生しない硬いものを何というか。}{}
\QQ{\RefBlank{sedimentary_rock}のうち、火山灰などの火山噴出物が積もってできたものを何というか。}{}
\QQ{生物の死がいや生活の跡が、重なった層の中に残ったものを\NamedBlank{fossil}という。}{}
\QQ{\RefBlank{fossil}のうち、層ができた当時の環境を知る手がかりとなるものを何というか。}{}
\QQ{\RefBlank{fossil}のうち、層ができた年代を知る手がかりとなるものを何というか。}{}
